% Source-bound extended proof of the selected paired cutoff comparisons.
\section{Sharp insertion and ordered cutoff estimates}
\label{pca:details}

Throughout this section, \(x\) denotes the erased squarefree cutoff
factor, \(I\) the positive insertion ideal, and \(B,K\) their respective
physical dyadic lengths. The exact partner norms are \(Nz=Nz_*Nx\) and
\(N\mathfrak k=Nk_*NI\). Here \(\mathfrak k\) is the original dual ideal,
and \(k_*\) is its fixed residual factor in a given cell. Write
\[
 E=Nd,\qquad A=A_1,\qquad R=A_1A_2,\qquad
 a=\log_D A,\quad r=\log_D R,\quad b=\log_D B,\quad k=\log_D K .
\]
At the critical pure scales \(a=r=3/4\) at the level of powers, and
at the critical mixed scales \(a=5/12,\ r=7/12\). Fixed comparison
constants in the scales are allowed throughout. A bound in this
section includes the common factor
\((DH)^\epsilon\mathcal M_{24}(W)^2\) for the paired cross, with the
original fixed-data constants. In row moments the appropriate one-profile
factor is understood separately.

The estimates concern the selected pairing with both original erasures,
all positive off-\(d\) increases whose sum with the old residue is at
most five, and the exact new residues on \(d\). The old cutoff and
conjugate erasures, including old exponent two in the conjugate
erasure, remain in the phase. The cutoff is \(Z=U^2\), with
\(U^2\le D\le7U^2/2\). The rows are the original unweighted literal
nonprincipal rows. We keep the common \(W\), fixed \(S,\nu,u_0\),
literal ray, capacity and squarefree selectors, and the actual
primitive \(Q_u\asymp R\), Gauss root and raw zeros. The physical
sieve is the order-six application in Section~\ref{phy:details}.
Its exponent at column length \(D^t\), relative to \(A\), is
\begin{equation}\label{pca:detail-delta}
 \Delta_a(t)=\max\bigl(0,t-a,t/3-a/6,5t/6-2a/3\bigr).
\end{equation}
Its growth slope is at most one. A live subunit dyad is treated
at constant scale.

\subsection{The middle subtraction and scalar comparisons}
\label{pca:detail-radial}

First remove the negligible inherited tail \(N\mathfrak k>D^5\).
Its absolute estimate applies to each selection below. Fix the
original allocation and dyad on the partner side. The compact
logarithmic separation of the partner profile uses
\(\log(NI/K)-\log(Nx/B)-\log(Q_u/R)\).
Its weighted Fourier absolute integral, with the extra factor
\(1+|\xi|\) in log frequency \(\xi\), is bounded by a constant
times \(\mathcal M_{24}(W)(1+y'_1)^{-4}\) at the actual
dyadic profile scale \(y'_1\). This follows by extending the
profile on a fixed-width log interval and integrating by parts.
The phase factors then separate on the erased array, insertion
array and row. There is no condition on \(B/K\) in this separation.
Sharp subinterval partial summation removes the insertion norm
twist with the erased array and every other row multiplier fixed.

For the original middle equality, the equation
\(Nz_*Nx=Nk_*NI\) fixes at most one integer norm on either side
given the other ideal. There are divisor-many ideals of each
integer norm on the clipped support. This counts the entire
insertion ideal, including its powerful part. The physical
\(1/Nx\) yields weighted mass
\(O((DH)^\epsilon\min(1,K/B))\). The two half radial powers
supply \((B/K)^{1/2}\), giving the factor
\(\sqrt{\min(B/K,K/B)}\le1\). The remaining label sums give
\(O(R/E)\). Apply this subtraction only in the original middle
window with its exact two endpoints. Equal norms in the
re-expressed direct, low and far pieces remain there.

The original live primal test gives
\begin{equation}\label{pca:detail-primal}
 E\ll D/B,\qquad
 D/(E E_1(Nh_1)^2)>Nz/4,\qquad y'_1\gg N\mathfrak k/R\gg K/R .
\end{equation}
Here \(h_1=f b_{\rm id}\) is the original partner label and
\(E_i=Ne_{1,i}\). We distinguish the ideal \(b_{\rm id}\) from
the length exponent \(b\). Let \(D^\rho\) be the product of the
two row norms relative to \(R\sqrt{K/B}\), after the projectors
and powerful sums. The full cutoff inverse norm is included
in these norms. For the exact paired scalar, start with
\[
 \frac{D}{E^2E_0E_1\,Ng\,N\ell\,Nz_*
          (Nf)^4(Nc)^2(Na_0)^2(Nb_{\rm id})^2}.
\]
Up to the fixed dyadic constants its profile scales are
\[
 \begin{aligned}
 y_0&=\frac{D}{EE_0 R(Nf\,Nc\,Na_0)^2},\\
 y'_1&\asymp
 \frac{D}{EE_1 R}\frac{(Na_0)^4}{(Nf\,Nc)^2}
 (N\ell_0 Nv_6)^6\frac{N\widetilde q_d}{Nj}
             \frac{Nk_S K}{B\,N\mathfrak e_\ell}.
 \end{aligned}
\]
The actual profile uses \(Q_u\), the physical ideal norms
and the original fixed \(S\)-factors.
Multiplying by \(y_0^{-\alpha}(y'_1)^{-\theta}\) gives the
following weights before the relative row moment:
\begin{equation}\label{pca:detail-scalar}
 \begin{aligned}
 &D^{1-\alpha-\theta}R^{\alpha+\theta}E^{\alpha+\theta-2}
   E_0^{\alpha-1} E_1^{\theta-1}(B/K)^\theta\\[-2pt]
 &\quad{}\times
 \frac{(Nf)^{2\alpha+2\theta-4}(Nc)^{2\alpha+2\theta-2}
       (Na_0)^{2\alpha-4\theta-2}}
      {Ng\,N\ell\,Nz_*(Nb_{\rm id})^2
        (N\ell_0 Nv_6)^{6\theta}(Nk_S)^\theta}
 (N\mathfrak e_\ell)^\theta
       (Nj/N\widetilde q_d)^\theta .
 \end{aligned}
\end{equation}
Here \(c\mid z_*,\bar\ell_{\rm g}\), \(a_0=a_{\rm id}/c\),
\(v_6\) is the full sixth-power repetition label,
\(\mathfrak e_\ell\) is the conjugate erasure with its actual
old exponents, and \(\widetilde q_d\) is the actual new representative
on \(d\). The scalar is accompanied by the retained decay
\((1+K/R)^{-(4-\theta)}\) from \eqref{pca:detail-primal}.
Every absolute comparison follows the exact pairing with those labels.

In the harmonic \(\ell\)-count the fixed \(\bar j\mid\ell\)
pays \(Nj\) to power one. A conjugate erasure prime \(p\) uses a
distinct \(\bar p\mid\ell\). Its conjugate lies in exactly one of
\(\ell_0,z_*,b_{\rm id}\), paying a second inverse norm. In the
old exponent-two case, the original \(p\mid z_*\) pays a third
inverse norm. Thus the resulting local exponents are at most
\(-2+\theta\) and \(-3+2\theta\), respectively. When \(Nc\)
or its half power is needed, charge \(c\mid z_*\). This charge
is disjoint from both conjugate charges since
\(\mathfrak q_\ell\perp a_{\rm id}\), its conjugates lie in
\(\ell\), and \(c\mid a_{\rm id},\bar\ell_{\rm g}\).
The other sums below are harmonic on polynomial support,
convergent inverse-power sums, or fixed-divisor choices.

For \(1/6\le\theta\le1\), take \(\alpha=1-\theta\).
The harmonic comparison at each \(d\) is
\begin{equation}\label{pca:detail-harmonic}
 \frac R E D^\rho (B/K)^{\theta-1/2}
                   (1+K/R)^{-(4-\theta)}.
\end{equation}
The choices \((\alpha,\theta)=(1,1/2)\) and \((7/6,1/3)\)
give, respectively, \(R^{3/2}D^\rho/(\sqrt D\sqrt E)\) and
\(R^{3/2}D^\rho(K/B)^{1/6}/(\sqrt D E^{1/3})\), with the
corresponding retained decays. In the second choice
\(E_0^{1/6}\ll_S E^{1/6}\). With \((\alpha,\theta)=(1,1/4)\),
the \(E_0\) cancels, the weights in \(f,a_0,\ell_0v_6\) have
exponents \(-3/2,-1,-3/2\), and the disjoint charge on \(c\)
pays its half power. The resulting comparison is
\begin{equation}\label{pca:detail-quarter}
 \frac{R^{5/4}D^\rho}{D^{1/4}E^{3/4}}
             (K/B)^{1/4}(1+K/R)^{-15/4}.
\end{equation}
By the first inequality in \eqref{pca:detail-primal}, summing
the \(E^{-3/4}\) costs \(O((D/B)^{1/4})\) on its live range.
The two three-half choices on \(E\le D^e\) cost
\(D^{e/2}\) and \(D^{2e/3}\), respectively.

The higher choice \((\alpha,\theta)=(5/2,1)\) leaves
\((Nf)^3(Nc)^5\) in \eqref{pca:detail-scalar}.
The exact primal restriction gives
\[
 (Nf)^2(Nb_{\rm id})^2Nz_*Nx\ll D/(E E_1),
 \qquad (Nf\,Nc)^4\ll (D/(E E_1B))^4 .
\]
Write the remaining numerator as
\((Nf)^{-1}Nc\,(Nf\,Nc)^4\) and pay the first two factors by
the preceding harmonic counts. The additional \(E^{-4}\) and
\(E_0\ll_S E\) give the harmonic comparison
\begin{equation}\label{pca:detail-high}
 \frac{R^{7/2}D^{3/2+\rho}}{B^4 E}
                 (B/K)^{1/2}(1+K/R)^{-3}.
\end{equation}

There is also a fixed-label selection for \((\alpha,\theta)=(1,0)\).
For \(p\mid a_0\), the original eligibility excludes \(b_{\rm id}\),
\(\ell,d,x,\mathfrak q_\ell\), and at most one of \(z_*,\bar\ell_{\rm g}\)
contains \(p\). The exact residual has valuation four or five at this
prime before adding the full sixth repetition. Thus
\(a_0^4v_6^6\mid k_*\). On \((Na_0)^4(Nv_6)^6\ll D^\xi\),
the joint count of \(a_0,v_6\) is \(O(D^{\xi/4})\), by summing
\((Nv_6)^{-3/2}\). Formula~\eqref{pca:detail-scalar} then gives
\(R D^{\rho+\xi/4}(K/B)^{1/2}/E\). This selection is sufficient
when \(Nk_*\ll D^\xi\). The comparison uses the full residual
ideal, with its \(S\)-part and repetitions.

\subsection{The sharp moment and the grouped sparse tail}
\label{pca:detail-sparse}

Write \(I=I_1P\) uniquely, where \(I_1\) is squarefree and the
nonzero valuations of \(P\) are \(2,3,4,5\). We have \(I_1\perp P\).
Fix \(P\) at first. The common-prime projector for \(x,I_1\)
uses \(s_1\mid(x,I_1)\), with \(\mu(s_1)\) and the two raw
values on \(s_1\). Require \(s_1\) itself to satisfy the original
fixed prime eligibility of \(I_1\), including \(s_1\perp P\),
before stripping it. The residual \(I'_1\) avoids \(s_1\)
by its squarefreeness, along with the original forbidden union.
Keep the original fixed \(P\)-deletion on the erased array.
After stripping \(s_1\), the squared erased mass is
\(O((DH)^\epsilon/(B Ns_1))\), including its actual \(1/Nx\).
The insertion length is \(K/(NP\,Ns_1)\).
A dyad \(NP\asymp M\) has \(O(\sqrt M)\) powerful ideals,
by the unique representation \(P=a_{\rm p}^2b_{\rm p}^3\)
with \(b_{\rm p}\) squarefree and the convergent
\(\sum_{b_{\rm p}}(Nb_{\rm p})^{-3/2}\). Overlap is allowed
at exponent five. This pays the \(M^{-1/2}\) shortening in
the product of the two moments.

For the sharp insertion bound at length \(L\), expand the
fixed deletion and squarefree selectors in the order
\(I'_1=e(r')^2m\), where \(e=e_f e_P\) divides the full forbidden
union and \(r'\) avoids that union. All Möbius signs and
fixed raw factors are retained. At fixed \(P,e,r'\) the remaining
sharp interval is an unweighted \(m\)-interval after the exact
partial summation. Its endpoints can be treated as sharp disk
endpoints. Smooth a disk of norm length \(L\) at radial width
proportional to \(\delta\sqrt L\). The normalized shell row
norm is bounded by a subpower times
\[
 D^{\Delta_a(l)/2}\bigl(\sqrt{\delta L}+L^{1/6}\bigr),
 \qquad l=\log_D L.
\]
To see the additive term, split each ideal uniquely as
\(\mathfrak s\,b_0 v^6\) with \(\mathfrak s\) supported on \(S\)
and \(b_0\) good sixth-free. At fixed \(\mathfrak s,v\),
the shell contains divisor-many ideals per integer norm and
has norm width \(O(\delta L/(N\mathfrak s(Nv)^6))\).
Apply the physical sieve to \(b_0\). The square roots of the
width weights sum absolutely. The constant atoms sum over
\(N\mathfrak s(Nv)^6\ll L\), giving \(L^{1/6}\) up to
subpowers. This also accounts for open or closed endpoints.

The smoothed primitive Poisson transform has the original
\(Q_u^{-1/2}\) and primitive Gauss root. On a dual dyad \(T\)
put \(Y=\sqrt{LT/R}\). The disk transform is \(L\) times
the Fourier transform of the unit disk's chord length
\(2\sqrt{1-v^2}\) at frequency comparable to \(Y\).
Integrating once by parts, separating the two endpoint
intervals of width \(1/Y\), and integrating by parts again
on the interior gives \(O(L(1+Y)^{-3/2})\).
The endpoint derivative integral is \(O(Y^{-1/2})\);
the interior boundary and second derivative terms
are \(O(Y^{1/2})\) before the second division by \(Y\).
The same argument with the chord length multiplied by \(v\)
bounds the frequency derivative, so its logarithmic derivative
costs one power of \(1+Y\). The convolution bump and its
log derivative provide \((1+\delta Y)^{-M_0}\)
for arbitrarily large fixed \(M_0\).
Consequently the cutoff log function \(m\) and \(m'\)
have \(L^2\) bounds \(O(L(1+Y)^{-3/2}(1+\delta Y)^{-M_0})\)
and \(O(L(1+Y)^{-1/2}(1+\delta Y)^{-M_0})\), respectively.
Splitting the Fourier integral at
\(1+\|m'\|_2/\|m\|_2\) and applying Cauchy and Plancherel
gives absolute log Fourier mass
\(O(L(1+Y)^{-1}(1+\delta Y)^{-M_0})\).
Multiplying by \(\sqrt T/\sqrt{Q_u}\) yields
\[
 O_{M_0}\!\left(\sqrt L
     (1+\delta\sqrt{LT/R})^{-M_0}\right).
\]
The zero frequency is absent by nonprincipality.
With \(\delta\asymp D^{\lambda-\Delta_a(l)}\)
when \(\Delta_a(l)>\lambda\), the effective dual exponent is
\begin{equation}\label{pca:detail-sharp-criterion}
 t=\max(0,r-l+2(\Delta_a(l)-\lambda)).
\end{equation}
The main shell term has row RMS \(O(\sqrt L D^{\lambda/2})\).
The additive term has the same bound whenever
\(\Delta_a(l)\le2l/3\). The smoothed term has that bound if
\(\Delta_a(t)\le\lambda\). Above the cutoff the arbitrary
smoothing decay pays the degree-one growth of the sieve and
the divisor subpowers. The ordinary estimate applies directly
if \(\Delta_a(l)\le\lambda\).

In pure use \(\lambda=0\) for \(l\le4/5\).
For \(l\ge3/8\) the convex function in
\eqref{pca:detail-sharp-criterion} before truncation is at
most \(3/8\) by the endpoints. In mixed use
\(\lambda=1/20\) for \(l\le5/8\). The smoothing branch
requires \(l>43/120\), and the same convex comparison gives
\(t\le43/120\). The additive term meets the stated criterion
in both cases. The exact deletion \(e\)-sum has divisor cost
at fixed \(P\), the \(r'\)-sum is harmonic, and the \(P\)-count
pays the shortening above. The resulting sharp rectangle
loss is \((\Delta_a(b)+\lambda)/2\). It uses the physical
ratio \(\sqrt{K/B}\) even when \(B\) and \(K\) are unequal.
The ordinary two-moment alternative has loss
\((\Delta_a(b)+\Delta_a(k))/2\).

The coupled comparison keeps a sharper tail for \(Nr'\ge D^q\)
at \(NP\asymp M\). Freeze \(e_f\), expand the remaining
\((x,P)=1\) by \(s_2\mid(x,P)\), with \(s_2\perp s_1\),
and fix the powered portion \(P_s\) on \(s_2\). Set \(P=P_sP'\)
and group the insertion coefficients by the actual
\(P'e_P(r')^2m\), of length \(K'=K/(Ns_1\,Ne_f\,NP_s)\).
The coefficients are divisor bounded and signed. They include
the original full \(P\)-dyad, threshold, old residue test and
clipped physical interval. For the sixth-free reduction of
this grouped ideal, write it as \(\mathfrak s\,b_0v^6\),
and split \(P'=P_0P_1,\ r'=r_0r_1\) by the support of \(v\).
The original \(r'\perp P\) gives \(P_1r_1^2\mid b_0\).
The two inverse-norm sums in the support count are bounded by
\[
 \min(1,(NP_sNP_0/M)^{1/2}),\qquad
 \min(1,Nr_0D^{-q}).
\]
The first is the powerful reciprocal dyad count and the
second the inverse-square ideal tail. The resulting squared
mass at fixed \(\mathfrak s,v,P_0,r_0\) is at most their
product times \(K'/(N\mathfrak s(Nv)^6)\), up to subpowers.
After square roots the \(v\)-sum converges, since
\((NP_0)^{1/4}(Nr_0)^{1/2}(Nv)^{-3}\le(Nv)^{-5/4}\).
The \(S\)-part sum also converges. Finally
\(NP_s\ge(Ns_2)^2\) makes the product of the insertion
shortening and the erased \(s_2\)-shortening harmonic.
Thus the relative norm saving in this branch is
\(M^{-1/4}D^{-q/2}\). It holds at the actual physical
lengths for both ratio orientations.

\subsection{Coupling the dual to the shorter ordered factor}
\label{pca:detail-coupling}

Open \(x\) into its actual two ordered cutoff factors.
Fix the assignments of the cutoff base and of \(s_1\), and
retain the individual cutoff bounds. Before expanding their
remaining coprimality, both factors are individually squarefree
and keep all their fixed deletions. Use the common squarefree
\(w\) projector, with pre-stripping eligibility in both arrays
and with both residual factors avoiding \(w\) after shortening.
The fixed phase is \(\mu(w)\psi_u(w)^2\). Sharp norm product
tests separate with polylogarithmic Fourier mass on the clipped
polynomial ranges, keeping their endpoints. The auxiliary
primitive insertion dual has the same raw orientation as \(x\).
Multiply it by the shorter cutoff factor as an actual ideal
before forming row norms. The grouped coefficients have
divisor bound and no added coprimality with the auxiliary
frequency. With shortened lengths \(X_s,X_l\) the physical
product is \(X_sX_l\asymp B/(Ns_1(Nw)^2)\).
The actual cutoff inverse norm and the smoothed transform
mass \(\sqrt L\) therefore give the factor
\[
 \frac{\sqrt{K/B}}{Ns_1 Nw\,Nr'\sqrt{Ne}\sqrt M}.
\]
The \(P\)-count pays the \(M^{-1/2}\), and the \(s_1,w,r'\)
sums remain harmonic.

For reference, set \(n=\log_D Ns_1,\ j=\log_D Ne,\ p=\log_D Nr'\),
\(v=\log_D M\), and write \(\ell=k-v-n-j-2p\) for the ordinary
insertion log length. In each coupled application choose the
relative target loss \(\rho\), then use the sharp branch for
large \(n\) and the extra deletion shortening \(D^{-j/2}\) for
large \(j\). Put
\[
 \sigma=\Delta_a(b)+\Delta_a(k)-2\rho,\qquad
 q=\max(0,\sigma-v/2).
\]
The grouped branch \(p\ge q\) pays the ordinary excess down to
\(\rho\). In the remaining branch take \(\delta\asymp D^{-\beta}\).
With
\(\beta=\max(0,\Delta_a(b-n)+\Delta_a(\ell)-j-2\rho)\)
the main shell meets the chosen loss directly.
The additive term has exponent
\((\Delta_a(b-n)+\Delta_a(\ell))/2-\ell/3-j/2\).
The smoothed dual exponent is
\[
 t=\max(0,r-k+v+n+j+2p+2\beta).
\]
Its two sieve losses are evaluated at \(u_s+t,u_l\), where
\(u_i=\log_D X_i\), \(u_s\le u_l\le1/2\), and \(u_s+u_l\le b-n\).
The degree bound pays all larger dual dyads using the same
smoothing decay. The following computations specify the
small-\(p\) shell and cutoff choices in the regions of interest.

\subsection{Pure central and comparable regions}
\label{pca:detail-central}

Set \(a=r=3/4\). On \(2/3\le b\le a,\ 1/2\le k\le31/40\)
the coupled loss is the smaller of \(\Delta_a(b)/2\) and
\[
 \rho_c=
 \begin{cases}
 (25b-2k)/204-7/272,& k\le a,\\
 (25b+22k)/204-31/272,& k\ge a .
 \end{cases}
\]
To verify the second choice when it is smaller, put
\(j_0=\Delta_a(b)-2\rho_c,\ n_0=3j_0\). Use the sharp branches
at \(n\ge n_0\) or \(j\ge j_0\), as described above. In the
remaining small-\(p\) range choose
\[
 \beta=\sigma-v/3-2n/3-4j/3-2p/3,\qquad
 6\rho_c=(a+b)/2-k+8\sigma/3 .
\]
Since \(v+2p\le2\sigma\), the lower bounds for \(\ell,\beta\)
on the two \(k\)-pieces are, respectively,
\[
 \begin{array}{c|cc}
 &\ell&\beta\\ \hline
 k\le a&47/68-9b/17+11k/51&73/408-9b/34+8k/153\\
 k\ge a&13/34-9b/17+32k/51&-97/204-9b/34+283k/306 .
 \end{array}
\]
Their minima are \(41/102,1/153\) and \(31/68,1/51\).
Thus \(\ell>3/8\) and \(\beta>0\).
The lower slope \(1/3\) gives
\(\Delta_a(\ell)\le\Delta_a(k)-(k-\ell)/3\);
the main shell costs at most \(\rho_c\).
For \(k\le a\) the additive excess is at most
\(-\ell/6-n/6-j/2\).
For \(k\ge a\) it is at most \(13/96-31/204<0\).

The dual satisfies \(t\le a-k+8\sigma/3\).
In the two \(k\)-pieces, \(b/2-k+8\sigma/3\) is
\((25b-2k-18)/34\) and \((25b+22k-36)/34\), respectively,
both negative on the stated rectangle. Thus \(u_s+t<a\).
The geometric mean excess is bounded by
\[
 \tfrac16\max(1/8,(a+b)/2-k+8\sigma/3)=\rho_c .
\]
Indeed both sieve arguments use the lower slope, with
\(u_s\le b/2-n/2,\ u_s+u_l\le b-n\) and \(u_l\le1/2\).

Write \(\rho\) for the smaller loss. Substituting it into the
harmonic comparison at \(\theta=1/6\) gives the short target
for \(57/80\le b\le a,\ k\le(172b+21)/264\).
The ordinary sharp loss already covers \(k\le b/2+3/16\).
At \(\theta=1\) and \(k\ge a\), the retained decay gives
the long target for \(273/404\le b\le a,\
k\ge(508b+1743)/2768\).
The sharp comparison covers \(k\ge5/8+4b/21\) through \(4/5\),
and the ordinary long insertion comparison decreases above \(4/5\).
The fixed product selection of Section~\ref{pca:detail-radial}
instead gives \(k\le(308b+21-204\xi)/400\) on the lower
\(k\)-piece. The two divisor-prefix choices with partner
exponents \(1/2,1/4\), retaining the actual partner decay,
give target through
\[
 e\le\max\left\{
 \tfrac14-2\rho+7(k-a)_+,\ 
 \tfrac14-4\rho+b-k+15(k-a)_+
 \right\},\qquad Nd\le D^e .
\]

For \(B\asymp K\) and \(b\le a\), the central formula gives
\(\rho\le1/17\), with the sharp branch sufficient below its
central rectangle. For comparable \(b\ge a\), the shifted
comparison below has loss at most
\(1/17+(1/51+5/12)(b-a)\) until the long target applies.
The partner decay in the half split pays this additional
loss. The two half powers and the one-plus-half powers then
give \(D^{55/68+\epsilon}\) over all \(d\) and the target
on \(Nd\le D^{9/68}\), respectively. The latter follows from
\(5/8+1/17+9/136=3/4\).

\subsection{Pure long insertion for every erasure length}
\label{pca:detail-long}

Take \(k=a+y,\ 0\le y\le1/56\). Work first at the reference
erased exponent \(a\), and put
\[
 \rho_0=1/17+y/51,\quad n_0=6(1/16-\rho_0),\quad
 j_0=n_0/3,\quad \sigma=9/68+27y/34 .
\]
The sharp and sparse branches use these reference thresholds,
with the actual insertion length. Their erased ordinary excess
increases by at most \(5(b-a)_+/6\), before its square root.
In the remaining branch use
\(\beta=\Delta_a(a-n)+\Delta_a(\ell)-j-2\rho_0\).
Here \(\ell\ge31/68-22y/51>3/8\) and
\(\beta\ge1/51-16y/153>0\).
The main shell has reference loss \(\rho_0\);
the additive term has excess at most
\(1/8+5y/12-(31/68-22y/51)/3<0\).

If \(\ell\le a\), then
\(\beta=\sigma-y/2-v/3-2n/3-4j/3-2p/3\).
If \(\ell\ge a\), then
\(\beta=\sigma-5v/6-7n/6-11j/6-5p/3\).
The first gives \(t\le8\sigma/3-2y\), and the second
\(t\le2\sigma-y\le8\sigma/3-2y=6\rho_0\).
At \(b>a\), reduce \(u_s\) to at most \((a-n)/2\) and
then \(u_l\) only enough to make the sum at most \(a-n\).
The reference pair remains ordered and nonnegative.
The total reduction is at most \(b-a\), either by the
shorter-factor cap alone or by the original bound on the sum.
The reference grouped length is below \(a\), since
\(3/8+6\rho_0<a\). Restoring the removed lengths adds at
most \(5(b-a)/12\) in the geometric mean, by the \(5/6\)
Lipschitz bound through actual exponent one.
Thus the relative loss is
\(\rho_0+5(b-a)_+/12\).

For \(b\le a\) use \eqref{pca:detail-harmonic} at \(\theta=1\);
for \(b\ge a\) use \eqref{pca:detail-quarter} with the
live divisor cap. The respective exponents after summation
are at most
\[
 a+\rho_0+(b-a)/2-7y/2,\qquad
 a+\rho_0-(b-a)/12-7y/2 .
\]
Both meet target at \(y\ge6/355\).
For \(k\ge a+1/56\) the sharp loss \(\Delta_a(b)/2\)
already suffices through \(4/5\), by the same two radial
comparisons. Above \(4/5\) the ordinary insertion adds
\(\Delta_a(k)/2\); at the endpoint
\(1/16+1/12-7/40<0\), and the bound decreases thereafter.
This gives \(K\gtrsim D^{1089/1420}\) over all \(B\) and \(d\).

\subsection{Pure long erasure and the intermediate gap}
\label{pca:detail-gap}

For \(a\le b\le1,\ k\le4/5\), the sharp loss is
\(\rho_{\rm sh}=5b/12-1/4\).
The three-half choice with partner exponent \(1/3\), summed
over \(E\ll D/B\), has exponent
\(a+5(1-b)/12-(a-k)_+/6-7(k-a)_+/2\).
The high radial comparison \eqref{pca:detail-high} has exponent
\(31/8-37b/12-k/2-3(k-a)_+\).
For \(k\le a\) the two reach target when
\[
 k\le L_b=-7/4+5b/2\quad\text{or}\quad
 k\ge H_b=25/4-37b/6 .
\]
They overlap for \(b\ge12/13\). For \(k\ge4/5\) the
ordinary insertion in the high comparison adds \(\Delta_a(k)/2\).
Its endpoint exponent \(409/120-37b/12\) is below target for
\(b\ge9/10\), and decreases with \(k\).
For the small residual product selection, the low endpoint
is instead \(1/2+b/6-\xi/2\); it overlaps the high endpoint
at \(b\ge69/76+3\xi/38\).

For \(1355/1482\le b\le12/13\), work in the gap
\(L_b\le k\le H_b\). Define
\[
 g=\min((k-L_b)/6,(H_b-k)/2),\qquad
 \rho=\rho_{\rm sh}-g,\qquad C_b=17/4-4b .
\]
The two losses agree at \(C_b\), where
\(\rho_C=3b/2-5/4\) is the minimum of \(\rho\).
It suffices to achieve the chosen \(\rho\), by using the
smaller radial comparison. Use
\(j_0=\Delta_a(b)-2\rho,\ n_0=6j_0/5,\ \sigma=j_0+\Delta_a(k)\),
and the exact nonnegative width exponent from
Section~\ref{pca:detail-coupling}. The two sharp branches and
the grouped tail apply with these choices. The minimum of
\(k-2\sigma-n_0-j_0\) occurs at \(C_b\) and equals
\((233b-202)/30>1/3\); this bounds \(\ell\) from below.
The shell additive excess is at most \(1/5-1/9<\rho_C\).

When \(\ell\ge3/8\) and \(\beta>0\), the slopes on \(b-n\)
and on \([\ell,k]\) bound \(t\) by
\(t_1=a-k+8\sigma/3\).
When \(\beta=0\), the bound is
\(t_2=a-k+2\sigma+n_0+j_0\).
For \(\ell<3/8,\ \beta>0\) use
\(\beta=j_0-5n/6-j\), which again gives \(t\le t_2\).
Both upper expressions peak at \(C_b\), as is seen from
the slopes \(1/3,-1\) of \(j_0\).
At \(C_b\) their difference is
\((438b-397)/180\ge0\), so \(t\le t_*=(95-96b)/18\).

For \(u_s+u_l\le b\), convexity and monotonicity of the
sieve losses reduce the geometric mean maximum to
\((u_s,u_l)=(b-1/2,1/2)\) and \((b/2,b/2)\).
Indeed the maximizing \(u_l\) at fixed \(u_s\) is
\(\min(1/2,b-u_s)\). Before the breakpoint use monotonicity
and after it use convexity. At the two endpoints the shifted
length is between \(a\) and one. The second endpoint is larger
by \((1-b)/8\), with value
\[
 7b/24+5t_*/12-5/16=815/432-139b/72\le\rho_C .
\]
The last inequality is exactly \(b\ge1355/1482\).
This closes the gap for every \(K\) after the two radial
comparisons and the independent middle subtraction.

\subsection{Pure short insertion over large erasures}
\label{pca:detail-short}

For \(b\ge a\) and \(k<a\), the harmonic choice \(\theta=1/6\)
and the quarter choice with the primal divisor cap meet target
if the rectangle loss is at most
\[
 \rho=\max\{(b-k)/3,\ b/2-k/4-3/16\}.
\]
With the sharp loss this covers \(b\le3-4k\) or
\(b\ge3k-3/4\), which is exhaustive for \(k\le15/28\).
For \(15/28\le k\le921/1652\), work inside the remaining
triangle and split it at \(b=9/8-k/2\).
Use the same \(j_0,n_0,\sigma\) and exact width as in the
intermediate gap. Here
\[
 j_0\le3/236,\quad b-n_0\ge3117/4130>a,\quad
 \ell\ge3159/8260>3/8,\quad \rho\ge1/16.
\]
These bounds follow by the affine formulas on the two triangles.
The main shell meets \(\rho\), and its additive term is at most
\((7/26+1/16)/2-1/8<1/16\).
The same \(t_1,t_2\) estimates apply, with
\(t_1-t_2\ge67/4130\). Thus it is enough to evaluate
\((\Delta_a(u_s+t_1)+\Delta_a(u_l))/2-\rho\)
at the two cutoff endpoints of the previous proof.
On each triangle these expressions are convex. Their
vertex values in the two endpoint orders are
\[
 \begin{array}{c|c|rr}
 k&b & (b-1/2,1/2)&(b/2,b/2)\\ \hline
 15/28&6/7&-5/168&-1/48\\
 921/1652&318/413&-83/9912&-83/9912\\
 921/1652&699/826&-51/3304&0\\
 921/1652&381/413&-787/19824&-85/2832 .
 \end{array}
\]
The sharp branches, sparse tail and smoothed dual therefore
meet \(\rho\) throughout. This proves the target for
\(K\lesssim D^{921/1652}\) and \(B\gtrsim R\) over all \(d\).

\subsection{Mixed ratios and the critical-capacity strip}
\label{pca:detail-mixed}

At \(a_*=5/12,\ r_*=7/12\) use the mixed sharp insertion loss
\(\lambda/2=1/40\) through \(k_0=5/8\). Denote
\(\Delta_{a_*}\) by \(\Delta\).
For \(k\ge b,\ k\le k_0\), the harmonic \(\theta=1\)
comparison is at most
\(r_*+\Delta(k)/2+1/40-3(k-r_*)_+\), maximal at \(k=r_*\).
For \(k\le b\le r_*\) use \(\theta=1/6\) and the sharp
moment, with the same bound \(57/80\). For \(k\le b,\ b\ge r_*\)
and \(k\le k_0\) the quarter comparison and live divisor cap
give exponent
\[
 23/48+\Delta(b)/2+1/40+(k-b)/4+(1-b)/4
                                      -15(k-r_*)_+/4 .
\]
At fixed \(b\) this is maximal at \(k=r_*\). It equals \(57/80\)
at \(b=r_*\), decreases with slope \(-1/12\) through \(5/6\),
and is constant thereafter. For \(k>k_0,\ k\ge b\) the ordinary
\(\theta=1\) exponent is at most \(r_*+\Delta(k)-3(k-r_*)\),
maximal at \(k_0\) with value \(101/144\). For \(k>k_0,\ k\le b\)
use \(\theta=1/6\). The exponent decreases in \(k\); at \(k=k_0\)
it is convex in \(b\) on \([k_0,1]\), with endpoint values
\(2/3,205/288\). Both are below \(57/80\).
Adding the separate harmonic equality count gives \(57/80\)
over all physical ratios and divisors.

Under \(A_1A_2^2\ll H\ll D^{3/4}\), the capacity bounds
\(r\le(3/4+a)/2\) at the level of exponents. The actual
\(Q_u\) is still used in each moment.
On the left, put \(a=a_*-z,\ 0\le z\le1/8\).
Use \(\lambda=1/20+3z/5\) through \(k_0=5/8\).
At \(l=43/120\) the three nonconstant monomials of
\(\Delta_a(l)\) are
\(-7/120+z,\ 1/20+z/6,\ 1/48+2z/3\), each at most \(\lambda\).
The smoothing branch is therefore longer than this \(l\).
The reference convex dual expression has endpoint maximum
\(31/90\), and the capacity and
\(\Delta_a(l)\le\Delta_{a_*}(l)+z\) give
\(t\le31/90+3z/10\). The three nonconstant monomials
at that upper bound are
\[
 -13/180+13z/10,\quad 49/1080+4z/15,\quad 1/108+11z/12 ,
\]
each at most \(\lambda\) by the endpoints in \(z\).
The additive term is controlled since \(3a\ge7/8\).
The erased and sharp insertion half-losses together increase
by at most \(4z/5\). In the harmonic branches the capacity
prefactor saves \(z/2\), while in the quarter branch it saves
\(5z/8\). In both ordinary branches \(k>5/8>r_*\), so the
active partner decay and capacity prefactor pay the increase
of the two sieve factors. This gives \(57/80+3z/10\le3/4\).

On the right take \(a=a_*+x,\ 0\le x\le59/912\), and use
\(\lambda=1/20+x/4\) through \(k_0=31/48\).
On the smoothing branch \(l>43/120\). In
\eqref{pca:detail-sharp-criterion} the possible \(r\)-increase
of \(x/2\) is paid by \(2(\lambda-1/20)\).
The reference convex function on \([43/120,31/48]\) has
endpoint values \(9/40,43/120\). Thus \(\Delta_a(t)\le1/20\);
the additive term is controlled by \(k_0<3a\).
In the ratio comparisons take the capacity upper bound
\(r_b=r_*+x/2\). The scalar estimates increase with \(r\),
while the sharp proof has already retained the actual conductor.
In the sharp ranges the maximum is
\(r_b+(\Delta_a(r_b)+\lambda)/2=57/80+x/2\).
For the quarter range, after setting \(k=r_b\),
the \(b\)-dependent term \((\Delta_a(b)-b)/2\) is
nonincreasing, giving the same maximum.
In the ordinary ranges the \(\theta=1\) endpoint is
\(21/32+4x/3\). The \(\theta=1/6\) bound at \(k=k_0\)
has \(b\)-endpoint values \(29/48+7x/4\) and
\(373/576+19x/12\); convexity bounds its interior.
All are at most \(3/4\), with the last endpoint attaining
that value at \(x=59/912\).
The union is the first-pivot strip
\(7/24\le a\le439/912\), subject to the original capacity
and with no independent closeness requirement on \(A_2\).

All comparisons here preserve the original \(\mu(d)\) in
the identity and use an absolute divisor bound at the end.
The remaining pure rectangles, the coherent divisor sum,
other signed components and the determinant and zero-location
comparisons remain open.
